% !PRINTED 233
\weberpage{233}
%<<CONT
of which the first is linear and the other is of the $n$th order, can be
determined in such a way that

\weq{7}{\begin{pmatrix} \alpha', & \beta' \\ \gamma', & \delta' \end{pmatrix}\begin{pmatrix} a, & 0 \\ c, & \partial \end{pmatrix} = \begin{pmatrix} 1, & 0 \\ 0, & n \end{pmatrix}\begin{pmatrix} \alpha, & \beta \\ \gamma, & \delta \end{pmatrix}.}

For if we write the relation (7) thus:
\[
\begin{pmatrix} \delta', & -\beta' \\ -\gamma', & \alpha' \end{pmatrix}\begin{pmatrix} 1, & 0 \\ 0, & n \end{pmatrix} = \begin{pmatrix} a, & 0 \\ c, & \partial \end{pmatrix}\begin{pmatrix} \delta, & -\beta \\ -\gamma, & \alpha \end{pmatrix},
\]
then one derives from it simply:

\weq{8}{\begin{aligned}
\delta' &= a\delta, & \gamma' &= \partial\gamma - c\delta,\\
\partial\beta' &= \beta, & n\alpha' &= \partial\alpha - c\beta, & a\alpha' &= \alpha - c\beta'.
\end{aligned}}

By this $\partial$ is first of all determined as the \gesperrt{greatest common
divisor} of $\beta$ and $n$, or [because of (4)] of $\mu'$ and $n$. For if
$\partial$ were not the greatest common divisor of $n=\partial a$ and
$\beta=\partial\beta'$, then $a$ and $\beta'$ would have to have a common
% [sic] The source prints \alpha=\alpha a'+c\beta, contrary to (8).
divisor, and consequently so would $\beta$ and $\alpha=\alpha a'+c\beta$, which
cannot be. At the same time $a=n/\partial$ is thereby determined, and $c$ must
satisfy the congruence

\weq{9}{\partial\alpha - c\beta \equiv 0\ (\mathrm{mod}\ n)}

by which, however, $c$ is determined not absolutely, but only to the modulus
$a$. One can therefore attach to $c$ the further condition, e.\,g., that it
shall be divisible by some power of 2, or, if $a$ is not divisible by 3, by
some power of 3, as we shall sometimes do later.

The three numbers $a$, $\partial$, $c$ can have no common divisor,
for such a one would [by (8)] also be a divisor of $\alpha$ and $\beta$, which
is impossible. If we therefore denote the greatest common divisor of $a$
and $\partial$ by $e$, then $c$ must be \gesperrt{relatively prime} to $e$.

The numbers $a$, $\partial$, $c$, the last modulo $a$, are, on account of (4)
and (8), completely determined by the two numbers $\mu$, $\mu'$, and do not
change when $\mu$, $\mu'$ are multiplied by a common factor coprime to $n$,
i.\,e. are replaced by another pair of the same \gesperrt{row}.
If one decomposes $n$ in any way into two factors $a$, $\partial$, then for
$c$ one can still take
\[
\frac{a}{e}\varphi(e)
\]
values to the modulus $a$ that are distinct and coprime to $e$. Each
of these assumptions leads to a pair of numbers $\mu$, $\mu'$ through the
congruences
%CONT>>
